\section{Formulación del problema}
\label{sec:formulacion_problema}

Con las dificultades encontradas (en 2025 no se podía recuperar la contraseña, no había donde registrar los documentos y no existía distribución por colegio), se plantea la siguiente pregunta principal para la gestión 2026:

\begin{center}
\textit{¿Cómo la página web de inscripción Versión 2.0 del Módulo Tecnológico Productivo San Julián BTH, en la gestión 2026, contribuye a resolver de forma completa la recuperación de la contraseña sin ir al colegio, el registro ordenado de los documentos y la distribución automática por colegio en cursos y paralelos, para que el proceso sea más justo, rápido y ordenado?}
\end{center}

De esta pregunta general surgen tres preguntas más específicas, relacionadas con cada problema:

\begin{itemize}
    \item ¿Cómo se puede lograr que el estudiante recupere solo y de forma segura su contraseña con su carnet y fecha de nacimiento, sin tener que hacer fila en el Módulo y asegurando que no pierda su cupo en la carrera que desea en 2026?
    \item ¿Cómo se puede organizar un espacio para que secretaría anote, revise y controle en el año que corresponde que se entregaron las libretas de 1ro a 3ro, la cédula de los tutores, el carnet y el certificado de nacimiento, con un registro ordenado y completo?
    \item ¿Cómo se puede lograr que el sistema distribuya de forma automática a los estudiantes en su curso y paralelo según el colegio del que provienen, para que los horarios coincidan y los cursos queden equilibrados?
\end{itemize}
